% Generated draft: requires mathematical and source review.
\documentclass[11pt]{article}
\usepackage[margin=25mm]{geometry}
\usepackage{fontspec}
\setmainfont{Latin Modern Roman}
\usepackage{amsmath,amssymb,mathtools,unicode-math}
\setmathfont{Latin Modern Math}
\usepackage{enumitem,needspace,xurl,hyperref}
\setlength{\parindent}{0pt}
\setlength{\parskip}{5pt}
\setlength{\emergencystretch}{3em}
\newcommand{\N}{\mathbb{N}}
\newcommand{\Z}{\mathbb{Z}}
\newcommand{\Q}{\mathbb{Q}}
\newcommand{\R}{\mathbb{R}}
\newcommand{\C}{\mathbb{C}}
\newcommand{\F}{\mathbb{F}}
\newcommand{\A}{\mathbb{A}}
\newcommand{\PP}{\mathbb{P}}
\newcommand{\E}{\mathbb{E}}
\newcommand{\eps}{\varepsilon}
\newcommand{\dd}{\,\mathrm{d}}
\begin{document}
\section*{UnsolvedMath ID: 30006456}
\textbf{Source problem code:} catalogue entry
\subsection{Definitions and mathematical statement}
For a nonzero integer \(d\), say that \(d\) is squarefree if no prime square divides \(d\), and write
\[
\mathcal S(N)=\{d\in\Z:0<|d|\leq N,\ d\text{ squarefree}\}.
\]
For a Dirichlet character \(\chi\), let \(L(s,\chi)\) denote the analytic continuation of the Dirichlet series
\[
L(s,\chi)=\sum_{n=1}^{\infty}\frac{\chi(n)}{n^s},\qquad \operatorname{Re}(s)>1.
\]
The source intends \(\chi_d\) to be a quadratic Dirichlet character associated with \(d\), but does not specify this association. Consequently, the following is the supplied conjectural assertion with an unresolved character convention, rather than a fully determined formulation.

For each fixed integer \(k\geq1\), there exists an explicit polynomial \(Q_k\), independent of \(N\), of degree exactly \(k(k+1)/2\), such that
\[
\frac{1}{N}\sum_{d\in\mathcal S(N)}L\!\left(\frac12,\chi_d\right)^k
=Q_k(\log N)+o(1)
\qquad(N\to\infty).
\]
Here \(\log\) denotes the natural logarithm, and \(o(1)\) means that, for each fixed \(k\), the difference between the displayed average and \(Q_k(\log N)\) tends to zero. No uniformity in \(k\) is asserted. The denominator is \(N\), as supplied, and both signs of \(d\) occur.

The character convention requires clarification: squarefree integers and fundamental discriminants are different indexing sets. For example, one possible convention associates a squarefree \(d\) with the fundamental discriminant
\[
D(d)=\begin{cases}d,&d\equiv1\pmod4,\\4d,&d\not\equiv1\pmod4,\end{cases}
\]
and then uses the Kronecker character \(\bigl(\frac{D(d)}{\cdot}\bigr)\). The supplied record does not establish that this is its intended convention. This draft therefore does not adopt that convention or replace the squarefree indexing condition by a fundamental-discriminant condition. The meaning of ``explicit'' is also left unspecified: neither coefficients of \(Q_k\) nor a precise recipe for them is supplied.
\subsection{Short English statement}
The supplied conjecture predicts a polynomial in \(\log N\), of degree \(k(k+1)/2\), for every fixed positive integral moment of central quadratic Dirichlet \(L\)-values, with an additive error tending to zero after division by \(N\). A precise statement still requires the intended definition of \(\chi_d\) for signed squarefree integers. This is a statement-only draft requiring human mathematical review.
\subsection{Sources}
\begin{itemize}
\item Supplied local UnsolvedMath registry record, permanent ID \(30006456\), title ``Moments of Quadratic Dirichlet L-Functions,'' exact source code \texttt{catalogue entry}. The original statement is labeled ``Conjecture 2.'' No individual proposer is supplied; the proposed year is recorded as 2025.
\item Supplied report citation: \emph{Analytic Number Theory (2026)}, Oberwolfach Reports, pp.\ 2701--2764, DOI \texttt{10.4171/owr/2025/51}. Supplied source and literature link: \url{https://doi.org/10.4171/owr/2025/51}, labeled ``Oberwolfach Report 51/2025.'' The difference between the title year and the report designation is retained without resolution.
\item Supplied literature reference: \emph{Integral moments of L-functions}, \url{https://doi.org/10.1112/S0024611505015174}. No authors or further bibliographic details are supplied in the record.
\item Literature assessment dated exactly \texttt{2026-08-22}: the supplied curation labels the problem ``partially solved'' and states that the full all-\(k\) conjecture remains open, with partial progress through low moments, bounds, and related function-field results. This is supplied dated curation derived from the OWR record, not an independently verified literature assessment. The source explicitly requests independent verification of time-sensitive status and cited claims.
\item The supplied background describes this quadratic-family conjecture as a special case of a general recipe for moments in families of \(L\)-functions. The recipe itself is not supplied. The registry label \texttt{corrected\_verified} is accompanied by a qualification that only formula delimiters, notation, grammar, and typography were repaired while preserving mathematical content and scope; it is not independent mathematical verification of this draft.
\end{itemize}
\end{document}
